\documentclass[11pt, UTF8]{ctexart}
\usepackage{amsmath, amsfonts, amssymb}
\usepackage{extarrows} 
\usepackage{geometry} 
\usepackage{xcolor} 
\usepackage{enumitem} 
\usepackage[most]{tcolorbox} 
\usepackage{fancyhdr} 
\usepackage{diagbox} % 用于绘制表格的斜线表头
\usepackage{tikz}    % 用于绘制概率密度积分区域图
\setlength{\headheight}{13.6pt}
\geometry{a4paper, margin=0.8in} 
\pagestyle{fancy}
\fancyhf{}
\lhead{\color{fblue} \hwdate 作业 答案} 
\rhead{\color{fblue} \today}
\cfoot{\thepage}

\definecolor{fblue}{RGB}{0, 74, 142}
\definecolor{fred}{RGB}{204, 26, 26}

\newtcolorbox{problem}[2]{
    enhanced,
    colback=white,
    colframe=fblue!80,
    arc=2mm,
    boxrule=0.8pt,
    before upper={\textbf{\color{fblue} 问题 #1 (#2)}\hspace{0.5em}},
    before skip=1.5em,
    after skip=1.5em,
    left=3mm, right=3mm, top=2mm, bottom=2mm
}

\newcommand{\hwdate}{4月2日}
\newcommand{\Prob}{P}
\newcommand{\comp}[1]{\overline{#1}\mskip 3mu}
\newcommand{\ans}[1]{\mathbf{\color{fred}#1}}

\begin{document}

\begin{center}
    {\LARGE \bfseries \color{fblue} \hwdate 作业 答案}
\end{center}

% ---------------- 题目 1 ----------------
\begin{problem}{1}{3.2}
(1) 盒子里装有 3 只黑球、2 只红球、2 只白球，在其中任取 4 只球. 以 $X$ 表示取到黑球的只数，以 $Y$ 表示取到红球的只数. 求 $X$ 和 $Y$ 的联合分布律. \\
(2) 在 (1) 中求 $P\{X > Y\}, P\{Y = 2X\}, P\{X + Y = 3\}, P\{X < 3 - Y\}$.
\end{problem}

\paragraph{\color{fblue}解} 
（1）按古典概型计算. 自 7 只球中取 4 只，共有 $\binom{7}{4} = 35$ 种取法. 在 4 只球中，黑球有 $i$ 只，红球有 $j$ 只（剩下 $4 - i - j$ 只为白球）的取法数为
\[ N\{X = i, Y = j\} = \binom{3}{i}\binom{2}{j}\binom{2}{4 - i - j}, \]
\[ i = 0,\, 1,\, 2,\, 3; \qquad j = 0,\, 1,\, 2; \qquad i + j \leqslant 4. \]
于是
\begin{align*}
    P\{X = 0, Y = 2\} &= \binom{3}{0}\binom{2}{2}\binom{2}{2}/35 = \frac{1}{35}. \\
    P\{X = 1, Y = 1\} &= \binom{3}{1}\binom{2}{1}\binom{2}{2}/35 = \frac{6}{35}. \\
    P\{X = 1, Y = 2\} &= \binom{3}{1}\binom{2}{2}\binom{2}{1}/35 = \frac{6}{35}. \\
    P\{X = 2, Y = 0\} &= \binom{3}{2}\binom{2}{0}\binom{2}{2}/35 = \frac{3}{35}. \\
    P\{X = 2, Y = 1\} &= \binom{3}{2}\binom{2}{1}\binom{2}{1}/35 = \frac{12}{35}. \\
    P\{X = 2, Y = 2\} &= \binom{3}{2}\binom{2}{2}\binom{2}{0}/35 = \frac{3}{35}. \\
    P\{X = 3, Y = 0\} &= \binom{3}{3}\binom{2}{0}\binom{2}{1}/35 = \frac{2}{35}. \\
    P\{X = 3, Y = 1\} &= \binom{3}{3}\binom{2}{1}\binom{2}{0}/35 = \frac{2}{35}.
\end{align*}
\[ P\{X = 0, Y = 0\} = P\{X = 0, Y = 1\} = P\{X = 1, Y = 0\} = P\{X = 3, Y = 2\} = 0. \]
分布律为
\begin{center}
\renewcommand{\arraystretch}{1.8}
\setlength{\tabcolsep}{20pt}
\begin{tabular}{c|cccc}
\hline
\diagbox{$Y$}{$X$} & $\displaystyle 0$ & $\displaystyle 1$ & $\displaystyle 2$ & $\displaystyle 3$ \\
\hline
$\displaystyle 0$ & $\displaystyle 0$ & $\displaystyle 0$ & $\displaystyle \frac{3}{35}$ & $\displaystyle \frac{2}{35}$ \\
$\displaystyle 1$ & $\displaystyle 0$ & $\displaystyle \frac{6}{35}$ & $\displaystyle \frac{12}{35}$ & $\displaystyle \frac{2}{35}$ \\
$\displaystyle 2$ & $\displaystyle \frac{1}{35}$ & $\displaystyle \frac{6}{35}$ & $\displaystyle \frac{3}{35}$ & $\displaystyle 0$ \\
\hline

\end{tabular}
\end{center}

（2）
\begin{align*}
P\{X > Y\} &= P\{X = 2, Y = 0\} + P\{X = 2, Y = 1\} \\
&\quad + P\{X = 3, Y = 0\} + P\{X = 3, Y = 1\} \\
&= \frac{3}{35} + \frac{12}{35} + \frac{2}{35} + \frac{2}{35} = \ans{\frac{19}{35}}. \\[1ex]
P\{Y = 2X\} &= P\{X = 1, Y = 2\} = \ans{\frac{6}{35}}. \\[1ex]
P\{X + Y = 3\} &= P\{X = 1, Y = 2\} + P\{X = 2, Y = 1\} + P\{X = 3, Y = 0\} \\
&= \frac{6}{35} + \frac{12}{35} + \frac{2}{35} = \ans{\frac{20}{35}} = \ans{\frac{4}{7}}. \\[1ex]
P\{X < 3 - Y\} &= P\{X + Y < 3\} \\
&= P\{X = 0, Y = 2\} + P\{X = 1, Y = 1\} + P\{X = 2, Y = 0\} \\
&= \frac{1}{35} + \frac{6}{35} + \frac{3}{35} = \ans{\frac{10}{35}} = \ans{\frac{2}{7}}.
\end{align*}

\vspace{1.5em}

% ---------------- 题目 2 ----------------
\begin{problem}{2}{3.3}
设随机变量 $(X, Y)$ 的概率密度为
\[ f(x, y) = \begin{cases} k(6 - x - y), & 0 < x < 2, 2 < y < 4, \\ 0, & \text{其他.} \end{cases} \]
(1) 确定常数 $k$. \\
(2) 求 $P\{X < 1, Y < 3\}$. \\
(3) 求 $P\{X < 1.5\}$. \\
(4) 求 $P\{X + Y \leqslant 4\}$.
\end{problem}

\paragraph{\color{fblue}解}
（1）由 $\int_{-\infty}^{\infty}\int_{-\infty}^{\infty} f(x, y)\mathrm{d}x\mathrm{d}y = 1$，得
\begin{align*}
1 &= \int_{2}^{4} \mathrm{d}y \int_{0}^{2} k(6 - x - y)\mathrm{d}x = k \int_{2}^{4} \left[ (6 - y)x - \frac{1}{2}x^2 \right]\Big|_{x=0}^{x=2} \mathrm{d}y \\
&= k \int_{2}^{4} (12 - 2y - 2)\mathrm{d}y = k(10y - y^2)\Big|_{2}^{4} = 8k,
\end{align*}
所以 $k = \ans{1/8}$.

（2）
\begin{align*}
P\{X < 1, Y < 3\} &= \int_{2}^{3} \mathrm{d}y \int_{0}^{1} \frac{1}{8}(6 - x - y)\mathrm{d}x \\
&= \frac{1}{8} \int_{2}^{3} \left[ (6 - y)x - \frac{1}{2}x^2 \right]\Big|_{x=0}^{x=1} \mathrm{d}y \\
&= \frac{1}{8} \int_{2}^{3} \left( \frac{11}{2} - y \right)\mathrm{d}y = \ans{\frac{3}{8}}.
\end{align*}

（3）
\begin{align*}
P\{X < 1.5\} &= \int_{2}^{4} \mathrm{d}y \int_{0}^{1.5} \frac{1}{8}(6 - x - y)\mathrm{d}x \\
&= \frac{1}{8} \int_{2}^{4} \left[ (6 - y)x - \frac{1}{2}x^2 \right]\Big|_{x=0}^{x=1.5} \mathrm{d}y \\
&= \frac{1}{8} \int_{2}^{4} \left( \frac{63}{8} - \frac{3}{2}y \right)\mathrm{d}y = \ans{\frac{27}{32}}.
\end{align*}

（4）在 $f(x, y) \neq 0$ 的区域 $R: 0 \leqslant x \leqslant 2, 2 \leqslant y \leqslant 4$ 上作直线 $x + y = 4$（如题 3.3 图），并记 $G: \{(x, y) \mid 0 \leqslant x \leqslant 2, 2 \leqslant y \leqslant 4 - x\}$，则

\noindent
\begin{minipage}{0.68\textwidth}
\begin{align*}
P\{X + Y \leqslant 4\} &= P\{(X, Y) \in G\} \\
&= \iint_{G} f(x, y)\mathrm{d}x\mathrm{d}y \\
&= \int_{2}^{4} \mathrm{d}y \int_{0}^{4-y} \frac{1}{8}(6 - x - y)\mathrm{d}x \\
&= \frac{1}{8} \int_{2}^{4} \left[ (6 - y)x - \frac{1}{2}x^2 \right]\Big|_{x=0}^{x=4-y} \mathrm{d}y \\
&= \frac{1}{8} \int_{2}^{4} \left[ (6 - y)(4 - y) - \frac{1}{2}(4 - y)^2 \right]\mathrm{d}y \\
&= \frac{1}{8} \int_{2}^{4} \left[ 2(4 - y) + \frac{1}{2}(4 - y)^2 \right]\mathrm{d}y \\
&= \frac{1}{8} \left[ -(4 - y)^2 - \frac{1}{6}(4 - y)^3 \right]\Big|_{2}^{4} = \ans{\frac{2}{3}}.
\end{align*}
\end{minipage}
\hfill
\begin{minipage}{0.28\textwidth}
\centering
\begin{tikzpicture}[scale=0.65, >=stealth]
\draw[->] (-0.5,0) -- (4.8,0) node[below] {$x$};
\draw[->] (0,-0.5) -- (0,4.8) node[left] {$y$};
\node[below left] at (0,0) {$O$};
\draw (2,0) -- (2,4);
\draw (0,2) -- (2,2);
\draw (0,4) -- (4,0);
\fill[black] (0,2) -- (2,2) -- (0,4) -- cycle;
\node[below] at (2,0) {$2$};
\node[below] at (4,0) {$4$};
\node[left] at (0,2) {$2$};
\node[left] at (0,4) {$4$};
\node[below] at (2,-1) {题 3.3 图};
\end{tikzpicture}
\end{minipage}

\vspace{1.5em}

% ---------------- 题目 3 ----------------
\begin{problem}{3}{3.7}
设二维随机变量 $(X, Y)$ 的概率密度为
\[ f(x, y) = \begin{cases} 4.8y(2 - x), & 0 \leqslant x \leqslant 1, 0 \leqslant y \leqslant x, \\ 0, & \text{其他.} \end{cases} \]
求边缘概率密度.
\end{problem}

\paragraph{\color{fblue}解}
$(X, Y)$ 的概率密度 $f(x, y)$ 在区域 $G: \{(x, y) \mid 0 \leqslant x \leqslant 1, 0 \leqslant y \leqslant x\}$ 外取零值. 如题 3.7 图有

\noindent
\begin{minipage}{0.65\textwidth}
\begin{align*}
f_X(x) &= \int_{-\infty}^{\infty} f(x, y)\mathrm{d}y \\
&= \begin{cases} \int_{0}^{x} 4.8y(2 - x)\mathrm{d}y, & 0 \leqslant x \leqslant 1, \\ 0, & \text{其他} \end{cases} \\
&= \begin{cases} \ans{2.4(2 - x)x^2,} & \ans{0 \leqslant x \leqslant 1,} \\ 0, & \text{其他.} \end{cases}
\end{align*}
\end{minipage}
\hfill
\begin{minipage}{0.3\textwidth}
\centering
\begin{tikzpicture}[scale=2, >=stealth]
\draw[->] (-0.2,0) -- (1.4,0) node[below] {$x$};
\draw[->] (0,-0.2) -- (0,1.4) node[left] {$y$};
\node[below left] at (0,0) {$O$};
\draw (0,0) -- (1.2,1.2) node[above left] {$y=x$};
\fill[black] (0,0) -- (1,0) -- (1,1) -- cycle;
\draw (1,0) -- (1,1);
\node[below] at (1,0) {$1$};
\node[below] at (0.5,-0.3) {题 3.7 图};
\end{tikzpicture}
\end{minipage}

\vspace{0.5em}

\begin{align*}
f_Y(y) &= \int_{-\infty}^{\infty} f(x, y)\mathrm{d}x \\
&= \begin{cases} \int_{y}^{1} 4.8y(2 - x)\mathrm{d}x, & 0 \leqslant y \leqslant 1, \\ 0, & \text{其他} \end{cases} \\
&= \begin{cases} \ans{2.4y(3 - 4y + y^2),} & \ans{0 \leqslant y \leqslant 1,} \\ 0, & \text{其他.} \end{cases}
\\
&= \begin{cases} \ans{2.4y(y-1)(y-3),} & \ans{0 \leqslant y \leqslant 1,} \\ 0, & \text{其他.} \end{cases}
\end{align*}

\textbf{注：}在求边缘概率密度时，需画出 $(X, Y)$ 的概率密度 $f(x, y) \neq 0$ 的区域，这对于正确写出所要求的积分的上下限是很有帮助的.

\end{document}